\begin{example}\label{ex:o2three}
三束等振幅、等相邻相位差光叠加，求主极大、零点与次极大，并画出一个周期内强度。
\Diagram{fig:o2threeq}{三个等长相量相邻夹角相同}{\draw[axis](0,0)--(2,0);\draw[axis](0,0)--(1.285575,1.532089);\draw[axis](0,0)--(-.347296,1.969616);\node[right]at(2,0){$A$};\node[above right]at(1.285575,1.532089){$A\ee^{i\delta}$};\node[left]at(-.347296,1.969616){$A\ee^{2i\delta}$};}
\end{example}
\begin{solution}
图\ref{fig:o2threeq}的三相量依次为$A,A\ee^{i\delta},A\ee^{2i\delta}$。提出公共相位得
\[
 \mathcal E=A\ee^{i\delta}(\ee^{-i\delta}+1+\ee^{i\delta})
 =A\ee^{i\delta}(1+2\cos\delta),
\]
所以$I/I_1=(1+2\cos\delta)^2$。零场条件为$\cos\delta=-1/2$，在$0\leq\delta\leq2\pi$内给$2\pi/3$、$4\pi/3$。
\Diagram{fig:o2threea}{三束干涉的两个零点和中央次极大}{\draw[axis](0,0)--(6.7,0)node[right]{$\delta$};\draw[axis](0,0)--(0,3.3)node[above]{$I/I_1$};\draw[curve,domain=0:6.283,samples=120]plot(\x,{(1+2*cos(\x r))^2/3});\foreach\x/\t in{2.094/{2\pi/3},3.142/\pi,4.189/{4\pi/3},6.283/{2\pi}}{\node[below]at(\x,0){$\t$};}\node[left]at(0,3){$9$};}
为区别极大与零点，对强度求导：
\[
 \frac1{I_1}\frac{\dd I}{\dd\delta}
 =-4(1+2\cos\delta)\sin\delta.
\]
除零场点外，驻点为$0,\pi,2\pi$。在$0,2\pi$三相量同向，强度为9；$\pi$处相量依次为$A,-A,A$，余下$A$，强度为1。由$\delta=\pi+\epsilon$时$1+2\cos\delta=-1+\epsilon^2+O(\epsilon^4)$，得到$I/I_1=1-2\epsilon^2+\cdots$，因此$\pi$确为次极大。

图\ref{fig:o2threea}的一个周期含一个次极大和两个完全暗点。周期平均强度为$\langle I\rangle/I_1=1+4\langle\cos\delta\rangle+4\langle\cos^2\delta\rangle=3$，等于三束单独强度之和；峰值9来自相干重排，不能解释为总能量增为三倍。
\end{solution}

\begin{example}\label{ex:o2unequal}
三束光的振幅依次为$A,2A,3A$，相邻相位差$\delta$。求强度、全局极值和可见度，判断是否有完全暗点。
\Diagram{fig:o2unequalq}{三个相量的长度为1、2、3，方向按$\delta$递增}{\draw[axis](0,0)--(1,0)node[right]{$1$};\draw[axis](0,0)--(1.285575,1.532089)node[right]{$2$};\draw[axis](0,0)--(-.520945,2.954423)node[left]{$3$};}
\end{example}
\begin{solution}
图\ref{fig:o2unequalq}给$\mathcal E/A=1+2\ee^{i\delta}+3\ee^{2i\delta}$。乘共轭时，每一对相量只产生一项余弦：
\begin{align*}
 I/I_1&=1+4+9+4\cos\delta+12\cos\delta+6\cos2\delta\\
 &=12\cos^2\delta+16\cos\delta+8.
\end{align*}
令$x=\cos\delta\in[-1,1]$，这是开口向上的二次式，顶点$x=-2/3$，最小值$8/3$。端点$x=1$强度36，$x=-1$强度4，所以全局最大36。
\Diagram{fig:o2unequala}{不等振幅的强度曲线不接触零轴}{\draw[axis](0,0)--(6.8,0)node[right]{$\delta$};\draw[axis](0,0)--(0,3.4)node[above]{$I/I_1$};\draw[curve,domain=0:6.283,samples=140]plot(\x,{(12*cos(\x r)^2+16*cos(\x r)+8)/12});\node[left]at(0,3){$36$};\draw[guide](0,.2222)--(6.3,.2222);\node[right]at(6.3,.35){$8/3$};}
图\ref{fig:o2unequala}没有完全暗点。全局可见度$(36-8/3)/(36+8/3)=25/29$。在$\delta=\pi$虽有局部次极大4，但它不决定全局可见度；求导时要区分关于$x$的端点和关于$\delta$的驻点。
\end{solution}

\begin{example}\label{ex:o2spacing}
三等振幅缝位置为$0,d,5d$。令$\delta=2\pi d\sin\theta/\lambda$，求强度和轴旁第一个零点。
\Diagram{fig:o2spacingq}{三缝的相邻间距为$d$和$4d$}{\draw[axis](-.3,0)--(6.1,0)node[right]{$x/d$};\foreach\x/\t in{0/0,1/1,5/5}{\draw[very thick](\x,-.4)--(\x,.5);\node[above]at(\x,.5){$\t$};}\draw[<->](1,-.8)--(5,-.8)node[midway,below]{$4d$};}
\end{example}
\begin{solution}
图\ref{fig:o2spacingq}给$\mathcal E/A=1+\ee^{i\delta}+\ee^{i5\delta}$，乘共轭得到$I/I_1=3+2\cos\delta+2\cos4\delta+2\cos5\delta$。完全相长在$\delta=2k\pi$，强度9。
\Diagram{fig:o2spacinga}{完全相消时三个等长相量构成正三角形}{\draw[->,thick](0,0)--(2,0)node[midway,below]{$1$};\draw[->,thick](2,0)--(1,1.732)node[midway,right]{$\ee^{i\delta}$};\draw[->,thick](1,1.732)--(0,0)node[midway,left]{$\ee^{i5\delta}$};}
三个等长矢量和为零必须彼此相差$2\pi/3$。由图\ref{fig:o2spacinga}，含有相位0和$\delta$的一对须满足$\delta=\pm2\pi/3$模$2\pi$；此时$5\delta$确为另一方向。因此首个正零点$\delta=2\pi/3$，$\sin\theta_0=\lambda/(3d)$。这一零点只有在$\lambda\leq3d$时属于传播角域；等号时位于掠射方向。不等间距不能直接把$N=3$代入等间距几何级数来解释全部旁瓣，但本题首零点恰好相同。
\end{solution}

\section{相干叠加的时间平均}
\subsection{交叉项由什么决定}
设两复包络为$\mathcal E_1(t),\mathcal E_2(t)$，共同载频为$\omega_0$。探测器先平均光学快振荡，再按曝光时间平均包络，强度含$\langle|\mathcal E_1|^2\rangle+\langle|\mathcal E_2|^2\rangle+2\operatorname{Re}\langle\mathcal E_1\mathcal E_2^*\rangle$。定义归一化互相干度
\[
 \gamma_{12}=\frac{\langle\mathcal E_1\mathcal E_2^*\rangle}
 {\sqrt{\langle|\mathcal E_1|^2\rangle\langle|\mathcal E_2|^2\rangle}},
 \quad |\gamma_{12}|\leq1.
\]
最后不等式由柯西不等式给出，保证总强度非负。相干度模控制对比度，其幅角控制条纹偏移；整体初相位虽可任取，两臂相对相位不能删去。

两个独立光源即使名义频率相同，相位也可能随机变化，使长曝光交叉项消失。稍异频而相位关系稳定的两束光，则可在短曝光中形成移动条纹；曝光足够长时，交叉项又会平均为零。是否记录到条纹，取决于频谱、相位变化与探测时间的共同作用。

\subsection{时间延迟与光程差}
同一稳定光源分成两路，一路相对另一路延迟$\tau$。取包含载频的复场$E(t)$，其互相关为$\langle E(t)E^*(t-\tau)\rangle$；不同频率分量互不相干时，这个量就是功率谱乘$\ee^{-i\omega\tau}$的频率积分。若使用前文除去载频的包络$\mathcal E(t)$，谱的变量应改为偏移频率$\Omega=\omega-\omega_0$，两种相关函数只差载频相位因子$\ee^{-i\omega_0\tau}$，模相同。矩形谱给sinc包络，双线谱给拍频包络，高斯谱给高斯包络。计算时先求每一频率的同频两路干涉，再将不同频率的强度相加。

在真空静态路径中，传播延迟为光程差除以$c$；色散材料中相位差为$\phi(\omega)$，一阶频率展开的群延迟是$\dd\phi/\dd\omega$。两臂在中心频率相位补偿为零，并不保证宽带光的群延迟也为零。白光干涉需要各频率的交叉相位在谱宽内不过度散开，因此一般要同时补偿中心相位和主要色散。

\subsection{空间延伸与相干面积}
扩展非相干光源的不同源元之间没有稳定相位关系，每个源元却可同时照亮两个观察点，产生同频相关项。设窄带光的波长为$\lambda$，源方向的近轴角坐标为$\boldsymbol\theta=(\theta_x,\theta_y)$，观察基线为$\vct d=(d_x,d_y)$。各源元的亮度权重为$B(\boldsymbol\theta)$，且两点到源的距离远大于源宽和基线。去掉与源元无关的共同相位，归一化互相干度为
\[
 \gamma(\vct d)=
 \frac{\iint B(\boldsymbol\theta)
 \ee^{ik(d_x\theta_x+d_y\theta_y)}\dd\theta_x\dd\theta_y}
 {\iint B(\boldsymbol\theta)\dd\theta_x\dd\theta_y},
 \qquad k=\frac{2\pi}{\lambda}.
\]
这里积分的是每个源元的两路相关项；独立源元之间的交叉项在平均后为零。它不是把独立源元的确定电场直接相加。

取图\ref{fig:o2rectcoherence}中的均匀矩形角源，两个方向的全角宽分别为$\Theta_x,\Theta_y$。积分可分离。对$i=x,y$，当$d_i\ne0$时，
\[
 \frac1{\Theta_i}\int_{-\Theta_i/2}^{\Theta_i/2}
 \ee^{ikd_i\theta_i}\dd\theta_i
 =\left[\frac{\ee^{ikd_i\theta_i}}{ikd_i\Theta_i}\right]_{-\Theta_i/2}^{\Theta_i/2}
 =\operatorname{sinc}\left(\frac{\pi\Theta_i d_i}{\lambda}\right);
\]
$d_i=0$时直接积分为1。因此
\[
 \gamma(d_x,d_y)=
 \operatorname{sinc}\left(\frac{\pi\Theta_xd_x}{\lambda}\right)
 \operatorname{sinc}\left(\frac{\pi\Theta_yd_y}{\lambda}\right).
\]
\Diagram{fig:o2rectcoherence}{矩形角源与基线平面上的首零范围；两个方向的宽度分别决定相干尺度}{
\begin{scope}
\fill[gray!15](-1.2,-.7)rectangle(1.2,.7);
\draw(-1.2,-.7)rectangle(1.2,.7);
\draw[axis](-1.9,0)--(2,0)node[right]{$\theta_x$};
\draw[axis](0,-1.3)--(0,1.8)node[above]{$\theta_y$};
\draw[<->](-1.2,-1)--(1.2,-1)node[midway,below]{$\Theta_x$};
\draw[<->](-1.5,-.7)--(-1.5,.7)node[midway,left,fill=white,inner sep=1pt]{$\Theta_y$};
\end{scope}
\begin{scope}[xshift=5.7cm]
\fill[gray!15](-.95,-.85)rectangle(.95,.85);
\draw[dashed](-.95,-.85)rectangle(.95,.85);
\draw[axis](-1.7,0)--(1.9,0)node[right]{$d_x$};
\draw[axis](0,-1.4)--(0,1.8)node[above]{$d_y$};
\node[above right]at(.95,0){$\lambda/\Theta_x$};
\node[above right]at(.1,.85){$\lambda/\Theta_y$};
\end{scope}}
两路等强时，可见度为$|\gamma|$。沿坐标轴的首零基线分别为$\lambda/\Theta_x$、$\lambda/\Theta_y$；首零围成的中央矩形面积为$4\lambda^2/(\Theta_x\Theta_y)$。若源实际宽度为$w_x,w_y$、距离为$a$，近轴有$\Theta_i\simeq w_i/a$，线源结果见例\ref{ex:o2extendedsource}。

以近轴相关函数的全平面延拓定义积分相干面积$A_c=\iint|\gamma(\vct d)|^2\dd d_x\dd d_y$。分别令$u_i=\pi\Theta_id_i/\lambda$，$\dd d_i=\lambda\dd u_i/(\pi\Theta_i)$，上下限仍为整个实轴，得到
\[
 A_c=\frac{\lambda^2}{\pi^2\Theta_x\Theta_y}
 \left(\int_{-\infty}^{\infty}\operatorname{sinc}^2u\,\dd u\right)^2
 =\frac{\lambda^2}{\Theta_x\Theta_y}.
\]
最后用了习题\ref{ex:o3energy}由帕塞瓦尔关系证明的积分$\int\operatorname{sinc}^2u\,\dd u=\pi$。这个积分面积与首零中央矩形面积采用不同定义；若按指定可见度阈值划区，系数还会改变。源角宽越小，相干基线范围越大；某方向角宽趋零时，该理想模型在对应基线方向上不衰减。

\section{薄膜程差的同波面比较}
膜内折射角为$r$，外部入射角为$i$，膜厚$h$。入射点$A$、底面反射点$B$、再出射点$C$的水平间距$AC=2h\tan r$。第二束在膜内累计光程$2nh/\cos r$，但两束要比较到同一个出射参考波面，第一束从$A$还走了一段$AD=AC\sin i$。若外部折射率$n_0$，相减得到
\begin{align*}
 \Delta&=\frac{2nh}{\cos r}-n_0(2h\tan r)\sin i\\
 &=\frac{2nh}{\cos r}-\frac{2nh\sin^2r}{\cos r}
 =2nh\cos r.
\end{align*}
第二步用$n_0\sin i=n\sin r$。扣除同一出射波面上的参考光程后，程差中便出现$\cos r$因子。

\Diagram{fig:o2filmphase}{$CD$垂直出射方向，$A\to D$光程须从膜内往返光程中扣除}{\draw[curve](-.7,0)--(3.2,0);\draw[curve](-.7,-1)--(3.2,-1);\fill(0,0)circle(1.3pt)node[left]{$A$};\fill(.57735,-1)circle(1.3pt)node[below]{$B$};\fill(1.1547,0)circle(1.3pt)node[below right]{$C$};\draw[->](0,0)--(.57735,-1)--(1.1547,0)--(2.8,1.6453);\draw[->](0,0)--(2,2);\draw[guide](.57735,.57735)--(1.1547,0);\fill(.57735,.57735)circle(1.3pt)node[above left]{$D$};}
图\ref{fig:o2filmphase}中两束出射光平行，$AD$沿第一束方向，$CD$为参考波面。若膜厚、折射率、入射角在观察区域变化，等相位条件相应形成等厚纹或等倾纹。实际观测透镜把哪一组光线汇到同一点，决定条纹的定域位置，不能只凭“是薄膜”就断言条纹一定在膜面。

总干涉相位还含反射相位。采用一致的偏振基时，两路实反射系数同号给相对相位0，异号给$\pi$；全反射的复系数则须取各自辐角。将它与传播相位相加后，才能确定亮暗级次。

\section{焦散相移的波动来由}
在腰面$z=0$取圆对称高斯场$U(r,0)=U_0\exp(-r^2/w_0^2)$。对于$z\ne0$，代入近轴菲涅耳传播积分，并展开$(x-\xi)^2+(y-\eta)^2$。记
\[
 a=\frac1{w_0^2}-\frac{ik}{2z},\qquad
 b_x=-\frac{ikx}{z},\qquad b_y=-\frac{iky}{z}.
\]
两个横向积分分别具有式\eqref{eq:o1gaussintegral}的形式，故
\begin{align*}
 U(r,z)&=\frac{U_0\ee^{ikz}}{i\lambda z}\ee^{ikr^2/(2z)}
 \mathcal G(a,b_x)\mathcal G(a,b_y)\\
 &=\frac{\pi U_0\ee^{ikz}}{i\lambda za}
 \exp\left[\frac{ikr^2}{2z}-\frac{k^2r^2}{4az^2}\right].
\end{align*}
取$z_R=kw_0^2/2$、$\lambda=2\pi/k$，记$\zeta=z/z_R$。则
\[
 aw_0^2=1-\frac{i}{\zeta},\qquad
 \frac{i\lambda za}{\pi}=1+i\zeta,\qquad
 \frac1{1-i/\zeta}=\frac{i\zeta}{1+i\zeta}.
\]
前因子因而为$U_0\ee^{ikz}/(1+i\zeta)$。指数中的两项先提取$r^2/w_0^2$，再通分：
\begin{align*}
 \frac{ikr^2}{2z}-\frac{k^2r^2}{4az^2}
 &=\frac{r^2}{w_0^2}
 \left[\frac i\zeta-\frac1{\zeta^2(1-i/\zeta)}\right]\\
 &=\frac{r^2}{w_0^2}
 \left[\frac i\zeta-\frac{i}{\zeta(1+i\zeta)}\right]
 =-\frac{r^2}{w_0^2(1+i\zeta)}.
\end{align*}
代回得到
\[
 U(r,z)=\frac{U_0\ee^{ikz}}{1+iz/z_R}
 \exp\left[-\frac{r^2}{w_0^2(1+iz/z_R)}\right].
\]
在轴上$r=0$，除去平面传播相位$kz$，剩余相位为$-\arctan(z/z_R)$。从负远处经过腰面到正远处，其总变化为$-\pi$。一维柱面高斯场只有一个高斯积分，前因子为$(1+iz/z_R)^{-1/2}$，额外相位总变化为$-\pi/2$。

\Diagram{fig:o2gouy}{高斯束腰附近，波动相位连续经过焦区}{\draw[axis](-3,0)--(3.2,0)node[right]{$z$};\draw[curve,domain=-2.7:2.7,samples=80]plot(\x,{.45*sqrt(1+\x*\x)});\draw[curve,domain=-2.7:2.7,samples=80]plot(\x,{-.45*sqrt(1+\x*\x)});\draw[guide](0,-1.6)--(0,1.6);\node[above]at(0,1.7){束腰};}
图\ref{fig:o2gouy}中的场在焦区有限、相位连续，几何光学常把整个焦区压成一个点，才把总变化记为一次焦点相移。焦点相移的数值取决于横向受限维数，因此点聚焦与线聚焦不同。两条干涉臂如果经历相同束腰和相同方向传播，这些相移通常相消，不应在每个透镜后重复添加。

\section{有限束数与多次反射腔的比较}
有限$N$束等幅的合场是有限几何级数，零点来自相量正多边形闭合；多次反射腔是无限几何级数，但每次场幅乘$R<1$，闭合多边形的零点结构不再相同。反射率越高，后续束衰减越慢，相干参与的有效往返数大致增多，透射共振更尖。

对正常入射腔，$2nl=m\lambda$，频率共振间隔$c/(2nl)$。若精细度为$\mathcal F_c$，线宽约$c/(2nl\mathcal F_c)$，中心频率$\nu_m=mc/(2nl)$，所以局部频率分辨本领$\nu_m/\delta\nu\simeq m\mathcal F_c$。它与有限光栅的$mN$形式类似，但$\mathcal F_c$由反射损耗和腔镜决定，不能把它严格解释成某个整数条数。

当腔镜不对称、存在吸收或散射，共振峰值一般不再为1，最佳入射耦合也依赖两镜透射率和内部损耗。前述单位峰值只属于理想无损对称腔。分析实测峰高时，应分别计入相位失配、孔径损失和吸收。

\subsection{扫描相位与扫描角度的极值区别}
对两束确定相干场，写$I(x)=I_1(x)+I_2(x)+2\sqrt{I_1(x)I_2(x)}|\gamma(x)|\cos\delta(x)$。若在固定空间点只改变一臂的均匀相位，前面的系数不变，最大、最小强度严格为$I_1+I_2\pm2\sqrt{I_1I_2}|\gamma|$，可见度公式可直接使用。

若沿屏幕扫描，包络、偏振投影和相干度也可能随$x$变化，实际极值应对整个$I(x)$求导。局部条纹可见度要求这些系数在一个纹距内近似不变。衍射包络零点、窄谱共振边缘或强渐变照明区需要保留系数的空间变化。可先在同一$x$固定位置扫描相位，区分包络变化与干涉对比度变化。
\Diagram{fig:o2localvisibility}{缓变包络内的局部明暗纹；靠近包络边缘时需对完整强度求极值}{\draw[axis](-.2,0)--(6.3,0)node[right]{$x$};\draw[axis](0,0)--(0,2.8)node[above]{$I$};\draw[guide,domain=0:5.8,samples=70]plot(\x,{2.4*exp(-.11*(\x-2.4)^2)});\draw[curve,domain=0:5.8,samples=200]plot(\x,{1.2*exp(-.11*(\x-2.4)^2)*(1+.8*cos(300*\x))});}
图\ref{fig:o2localvisibility}的上包络不等于每个实际亮纹都恰好落在$\delta=2m\pi$；相位极值与包络斜率共同决定峰位。这一点在有限$N$束加权、法布里—珀罗有吸收以及不等宽双缝中都会出现。

多束完全相干时还应区别场相加与功率归一化。若场权重为$a_j$，主峰场为$\sum a_j$，而输入功率按$\sum|a_j|^2$计。改变权重后比较峰高，必须保持同一功率条件；否则峰高的改变可能仅来自总照明功率改变。旁瓣降低、主峰变宽和峰值变化应分别从同一个加权和推导，不能用其中一个指标代替全部光学性能。

\section{有限带宽、有限光源与曝光时间}
\begin{example}\label{ex:o2bandwidth}
双束干涉的功率谱在频率$\nu_0-\Delta\nu/2$到$\nu_0+\Delta\nu/2$内恒定，其他处为零。不同频率互不相干，同频两路完全相干且等强。两臂时间差为$\tau$，求干涉可见度包络及首个零点。
\Diagram{fig:o2bandq}{矩形功率谱的每个频率分别产生自己的干涉项}{\draw[axis](0,0)--(5,0)node[right]{$\nu$};\draw[axis](0,0)--(0,2.8)node[above]{功率谱};\draw[curve](1.2,0)--(1.2,2)--(3.8,2)--(3.8,0);\node[below]at(2.5,0){$\nu_0$};\node[above]at(2.5,2){$\Delta\nu$};}
\end{example}
\begin{solution}
图\ref{fig:o2bandq}总干涉项需对频率积分。先计算长度为$L$的均匀区间平均：当$b\ne0$时，
\begin{align}
 \frac1L\int_{-L/2}^{L/2}\cos(a+bs)\dd s
 &=\left[\frac{\sin(a+bs)}{bL}\right]_{-L/2}^{L/2}\notag\\
 &=\frac{2\cos a\sin(bL/2)}{bL}
 =\cos a\operatorname{sinc}(bL/2).\label{eq:o2rectaverage}
\end{align}
当$b=0$时直接积分得$\cos a$，与连续极限相同。对正弦作同一计算，得$\sin a\operatorname{sinc}(bL/2)$，因此复指数的平均为$\ee^{ia}\operatorname{sinc}(bL/2)$。本题以$\nu=\nu_0+s$换元，取$L=\Delta\nu,a=2\pi\nu_0\tau,b=2\pi\tau$，得到归一化干涉项$\cos(2\pi\nu_0\tau)\operatorname{sinc}(\pi\Delta\nu\tau)$。
\Diagram{fig:o2banda}{可见度是谱积分的归一化模，有限带宽使远离零光程差的条纹变弱}{\draw[axis](-3.2,0)--(3.4,0)node[right]{$\Delta\nu\tau$};\draw[axis](0,0)--(0,2.7)node[above]{$\mathcal V$};\draw[curve,domain=.01:3,samples=80]plot(\x,{2.2*abs(sin(180*\x)/(pi*\x))});\draw[curve,domain=-3:-.01,samples=80]plot(\x,{2.2*abs(sin(180*\x)/(pi*\x))});}
故图\ref{fig:o2banda}的可见度$\mathcal V=|\operatorname{sinc}(\pi\Delta\nu\tau)|$，首零$|\tau|=1/\Delta\nu$，对应光程差尺度$c/\Delta\nu$。这一定义依赖所选的矩形谱及“首零”标准；高斯谱没有有限首零，不能把相干长度当所有光源共有的硬截止距离。
\end{solution}

\begin{example}\label{ex:o2extendedsource}
宽$w$的均匀线光源平行于双缝连线，位于缝前距离$a$。源上不同点互不相干，每个点各自照明两缝。双缝距$d$。求近轴干涉可见度及第一消失源宽。
\Diagram{fig:o2sourceq}{扩展光源各点产生相位不同的条纹，相加强度}{\draw[very thick](0,-1.3)--(0,1.3);\node[left]at(0,0){$w$};\fill(3,.6)circle(1.5pt)node[right]{$S_+$};\fill(3,-.6)circle(1.5pt)node[right]{$S_-$};\draw[guide](0,1)--(3,.6);\draw[guide](0,1)--(3,-.6);\node[below]at(1.5,-1.5){$a$};}
\end{example}
\begin{solution}
图\ref{fig:o2sourceq}横坐标$x_s$的源元附加双缝相位$2\pi dx_s/(\lambda a)$。屏上该点的共同相位记为$\delta_0$。在式\eqref{eq:o2rectaverage}已展开的定积分中，取区间长度$w$、相位斜率$2\pi d/(\lambda a)$，均匀源强度平均的交叉项比例为
\[
 \frac1w\int_{-w/2}^{w/2}\cos\left(\delta_0+\frac{2\pi dx_s}{\lambda a}\right)\dd x_s
 =\cos\delta_0\,\operatorname{sinc}\left(\frac{\pi dw}{\lambda a}\right).
\]
\Diagram{fig:o2sourcea}{源角宽越大，不同源点的中央纹偏移越难重合}{\draw[axis](0,0)--(5.3,0)node[right]{屏坐标};\draw[curve,domain=0:4.7,samples=100]plot(\x,{1+.5*cos(300*\x)});\draw[guide,domain=0:4.7,samples=100]plot(\x,{1+.5*cos(300*\x+120)});\node[above]at(2.5,2.1){不同源点条纹示意};}
所以$\mathcal V=|\operatorname{sinc}(\pi dw/(\lambda a))|$，首零$w=\lambda a/d$。图\ref{fig:o2sourcea}表达的是不同条纹强度相加，不能先把独立源点的电场相干相加。源尺寸沿另一方向延伸时，其投影到双缝基线的宽度才进入本式。
\end{solution}

\begin{example}\label{ex:o2movingfringes}
两平面波频率分别为$\omega_1,\omega_2$，波矢分别$\vct k_1,\vct k_2$，偏振相同。求等相位差条纹的运动速度，并讨论曝光时间$T$对可见度的影响。
\Diagram{fig:o2movingq}{两个不同频率和方向的平面波形成移动干涉纹}{\draw[axis](0,0)--(3.2,1.2)node[above]{$\vct k_1$};\draw[axis](0,0)--(3.2,-1.2)node[below]{$\vct k_2$};\foreach\y in{-1,0,1}{\draw[curve](4,\y)--(6,\y);}\draw[axis](6.3,-1)--(6.3,1)node[right]{$\Delta\vct k$};}
\end{example}
\begin{solution}
图\ref{fig:o2movingq}交叉相位为$\Delta\vct k\cdot\vct r-\Delta\omega t+\phi$。沿条纹随动点求时间导数，$\Delta\vct k\cdot\vct v_f-\Delta\omega=0$。取速度沿条纹法线，且$\Delta\vct k\ne0$，得到$\vct v_f=(\Delta\omega/|\Delta\vct k|^2)\Delta\vct k$。若$\Delta\vct k=0$，干涉项在空间中均匀；$\Delta\omega\ne0$时只有时间拍动，二者同时为零时则相位差恒定，均无可定义的空间条纹移动速度。纹速描述等相位面的移动，光能和信息仍按各束光的传播规律传递。
\Diagram{fig:o2movinga}{固定像素在有限曝光内平均正弦变化，长曝光削弱干涉项}{\draw[axis](0,0)--(5.5,0)node[right]{$t$};\draw[axis](0,0)--(0,2.7)node[above]{光强};\draw[curve,domain=0:5,samples=130]plot(\x,{1.3+cos(360*\x)});\draw[guide](1,0)--(1,2.5);\draw[guide](3.5,0)--(3.5,2.5);\node[above]at(2.25,2.5){曝光$T$};}
取曝光变量$s=t-t_0$，则交叉相位为$a-\Delta\omega s$，其中$a=\Delta\vct k\cdot\vct r-\Delta\omega t_0+\phi$。在式\eqref{eq:o2rectaverage}中取$L=T,b=-\Delta\omega$，交叉项便乘$\operatorname{sinc}(\Delta\omega T/2)$，等强可见度为其绝对值。图\ref{fig:o2movinga}在$T\gg1/|\Delta\omega|$时条纹变弱，短曝光则可记录移动条纹；“不同频率不干涉”包含了长时间平均的条件。
\end{solution}

\section{多次反射与干涉测量}
\begin{example}\label{ex:o2fabry}
无损对称法布里—珀罗平行腔两镜的能量反射率均为$R$，腔内单程相位$\phi$。取镜的振幅系数使腔往返倍率为$R\ee^{2i\phi}$。求透射强度并推导高反射率时的精细度。
\Diagram{fig:o2fabryq}{腔内每多一次往返，场乘$R\ee^{2i\phi}$}{\draw[very thick](1,-1.3)--(1,1.3);\draw[very thick](4,-1.3)--(4,1.3);\draw[->](0,.8)--(4,.8);\draw[->](4,.5)--(1,.5);\draw[->](1,.2)--(4,.2);\draw[->](4,-.1)--(5.5,-.1);\node[below]at(2.5,-1.4){$\phi$};}
\end{example}
\begin{solution}
图\ref{fig:o2fabryq}中，每个无损镜的能量透射率为$1-R$，首束透射场模为$(1-R)A$，包含单程相位$\ee^{i\phi}$。后续每束再乘一次往返倍率$R\ee^{2i\phi}$，几何级数给
\[
 A_t=(1-R)A\ee^{i\phi}\sum_{j=0}^\infty R^j\ee^{2ij\phi}
 =\frac{(1-R)A\ee^{i\phi}}{1-R\ee^{2i\phi}}.
\]
分母模平方为$1+R^2-2R\cos2\phi=(1-R)^2+4R\sin^2\phi$，因此
\[
 \frac{I_t}{I_i}=\frac1{1+\mathcal A\sin^2\phi},\qquad
 \mathcal A=\frac{4R}{(1-R)^2}.
\]
\Diagram{fig:o2fabrya}{高反射率使透射峰变窄，理想共振峰值仍为1}{\draw[axis](0,0)--(6.7,0)node[right]{$\phi$};\draw[axis](0,0)--(0,2.7)node[above]{$I_t/I_i$};\draw[curve,domain=0:6.283,samples=220]plot(\x,{2.3/(1+80*sin(\x r)^2)});\node[below]at(3.142,0){$\pi$};}
图\ref{fig:o2fabrya}的峰在$\phi=m\pi$，峰高为1。若$\mathcal A\geq1$，峰两侧半高点满足$\sin^2\delta\phi=1/\mathcal A$，故准确半高全宽为$2\arcsin(1/\sqrt{\mathcal A})$，精细度为
\[
 \mathcal F_c=\frac{\pi}{2\arcsin[(1-R)/(2\sqrt R)]}.
\]
高反射率时$(1-R)/(2\sqrt R)\ll1$，用$\arcsin x\simeq x$得到$\mathcal F_c\simeq\pi\sqrt R/(1-R)$。若$\mathcal A<1$，整条曲线的最低值仍高于峰的一半，按这个基线为零的半高定义没有半高点，不应使用窄峰线宽近似。腔内各束振幅逐次衰减，属于无限几何级数。无损对称腔在共振处可完全透射，与有限等幅多束的$N^2$峰高采用不同归一化。
\end{solution}

\begin{example}\label{ex:o2gascell}
迈克耳孙一臂放入长度$l$的气室，光在其中往返。抽真空后充入折射率$n$的气体，记录中心通过$N$条纹。求$n$；若只改变另一臂镜子的位置补偿，需移动多少？
\Diagram{fig:o2gasq}{气室在同一干涉臂内被光往返通过}{\draw[axis](0,0)--(6.3,0);\draw[curve](2,-.8)rectangle(4,.8);\draw[very thick](5.5,-1.2)--(5.5,1.2);\node[above]at(3,.8){气室$n,l$};\draw[->](1,.3)--(5.3,.3);\draw[->](5.3,-.3)--(1,-.3);}
\end{example}
\begin{solution}
图\ref{fig:o2gasq}以抽空状态为参考，每一程附加光程为$(n-1)l$，往返总改变量$2(n-1)l=N\lambda_0$，所以$n=1+N\lambda_0/(2l)$。
\Diagram{fig:o2gasa}{另一臂镜移一段，往返补偿两倍该长度}{\draw[very thick](2,-1)--(2,1);\draw[guide](3.3,-1)--(3.3,1);\draw[<->](2,-1.5)--(3.3,-1.5)node[midway,below]{$\Delta d$};\node[right]at(3.6,0){$2\Delta d=2(n-1)l$};}
图\ref{fig:o2gasa}需要镜移大小$|\Delta d|=(n-1)l$，方向选为增加另一臂光程来补偿气室臂的增加。若气室窗在充气前后保持不变，窗片的静态相位相消；若充气使窗形变或气室长度变化，测得的不再单是折射率项，需把额外光程纳入。
\end{solution}

\begin{example}\label{ex:o2ringmotion}
牛顿环中心原间隙$h_0$，将透镜整体抬高$\Delta h>0$且不改变曲率。跟踪固定级次$m$的反射暗环，求半径平方变化和中心每暗一次所需的抬升量。
\Diagram{fig:o2ringmotionq}{整体抬高改变所有半径处的共同间隙}{\draw[curve](-2.3,0)--(2.3,0);\draw[curve,domain=-2:2,samples=60]plot(\x,{.2+.15*\x*\x});\draw[guide,domain=-2:2,samples=60]plot(\x,{.7+.15*\x*\x});\draw[->](0,.25)--(0,.65)node[right]{$\Delta h$};}
\end{example}
\begin{solution}
图\ref{fig:o2ringmotionq}的暗环$r_m^2=R[m\lambda_0/n-2h_0]$，抬高后$r_m'^2=r_m^2-2R\Delta h$。所以同一级次向内收缩，直到$r_m'=0$后消失。
\Diagram{fig:o2ringmotiona}{跟踪同一级次时圆环向中心收缩}{\draw[curve](0,0)circle(1.6);\draw[guide](0,0)circle(1.05);\draw[->](1.55,0)--(1.1,0);\draw[->](-1.55,0)--(-1.1,0);\fill(0,0)circle(1pt);}
中心相位的暗纹条件$2nh_0=m\lambda_0$，相邻暗态间隔$\Delta h=\lambda_0/(2n)$。图\ref{fig:o2ringmotiona}所说向内是固定$m$的跟踪；如果每次都把“从中心数的第一环”重新编号，编号跳变会掩盖运动方向，必须先说明跟踪对象。
\end{solution}

\begin{example}\label{ex:o2weightedbinomial}
四束等间距光的振幅依次为$1,3,3,1$倍$A$。求干涉强度，比较它与四束等振幅时的零点和旁瓣，讨论权重改变的代价。
\Diagram{fig:o2binomialq}{四束相邻相位差为$\delta$，振幅取二项式权重}{\draw[axis](-.3,0)--(4.3,0)node[right]{束序};\foreach\x/\y in{0/1,1/3,2/3,3/1}{\draw[very thick](\x,0)--(\x,{.7*\y});\node[above]at(\x,{.7*\y}){$\y$};}}
\end{example}
\begin{solution}
图\ref{fig:o2binomialq}复振幅$\mathcal E/A=1+3\ee^{i\delta}+3\ee^{2i\delta}+\ee^{3i\delta}=(1+\ee^{i\delta})^3$，所以$I/I_A=64\cos^6(\delta/2)$。
\Diagram{fig:o2binomiala}{二项式加权消去周期内旁瓣，但主峰更宽}{\draw[axis](0,0)--(6.8,0)node[right]{$\delta$};\draw[axis](0,0)--(0,2.7)node[above]{归一化强度};\draw[curve,domain=0:6.283,samples=100]plot(\x,{2.3*cos((\x/2) r)^6});\node[below]at(3.142,0){$\pi$};}
图\ref{fig:o2binomiala}一个周期内只有$\delta=\pi$的零点，没有次极大。四束等幅时有三个零点和两个次极大、主峰首零在$\pi/2$。加权减少旁瓣却加宽主峰；若总入射功率固定，还应把权重整体归一化，不能直接拿峰值64与等幅16比较“效率”。
\end{solution}

\section{腔内损耗、耦合匹配与相位测量}
\subsection{把振幅损耗与能量损耗分别保留}
腔的透射峰取决于两端耦合和内部损耗。设两镜能量反射率为$0\leq R_1,R_2\leq1$，镜本身无吸收，腔内单程的能量存活率为$0\leq\eta\leq1$。一次往返的振幅模为$q=\eta\sqrt{R_1R_2}$，以下稳态级数要求$q<1$。用$\delta$表示包括镜相移在内的往返相位，略去各束共有的单程相位后，透射场为
\[
 E_t=\sqrt{(1-R_1)(1-R_2)\eta}\,E_{
 \rm in}\sum_{j=0}^{\infty}q^j\ee^{\mathrm i j\delta}.
\]
因此
\[
 T(\delta)=\frac{(1-R_1)(1-R_2)\eta}
 {(1-q)^2+4q\sin^2(\delta/2)}.
\]
此处$\sqrt\eta$是单程振幅因子，$\eta$才是往返振幅因子。两者混用会使峰高和线宽同时错误。

在$\delta=2m\pi$附近取$\delta=2m\pi+\epsilon$，若$q$接近1，则半高点满足$q\epsilon^2=(1-q)^2$，故往返相位全宽为$2(1-q)/\sqrt q$。精确半高点要求$\sin^2(\epsilon/2)=(1-q)^2/(4q)$；右边超过1时，一个自由光谱区内没有下降到峰值一半的位置，便不宜继续套用高精细度线宽公式。

\ExpansionDiagram{fig:dep-o2-roundtrip}{每次新增透射束比前一束多一次往返及其损耗}{
\draw[very thick](0,-1)--(0,1);\draw[very thick](4,-1)--(4,1);\node[above]at(0,1){$R_1$};\node[above]at(4,1){$R_2$};
\draw[->](-1,.6)--(4,.6);\draw[->](4,.2)--(0,.2);\draw[->](0,-.2)--(5,-.2);\node[below]at(2,-.6){单程能量因子$\eta$};}
图\ref{fig:dep-o2-roundtrip}中的一次往返含两次单程衰减，但振幅每次只乘$\sqrt\eta$。把它们与两镜反射振幅相乘，才得到$q$。

\begin{example}有内损耗的腔怎样消除入射端反射
采用上述模型，入射端的振幅反射系数取
\[
 r_{\rm cav}(\delta)=\frac{\sqrt{R_1}-\eta\sqrt{R_2}\ee^{\mathrm i\delta}}
 {1-\eta\sqrt{R_1R_2}\ee^{\mathrm i\delta}}.
\]
求共振时反射为零的条件、此时透射率与腔内吸收率。解释反射为零是否意味着全部透射。
\ExpansionDiagram{fig:dep-o2-match-q}{入射端反射由直接反射和腔内返回光相干组成}{
\draw[very thick](1,-1)--(1,1);\draw[very thick](5,-1)--(5,1);
\draw[->](-1,.5)--(1,.5);\draw[->](1,.1)--(-1,.1)node[left]{直接反射};\draw[->](5,-.3)--(-1,-.3)node[left]{返回光};\node[above]at(3,1){有损腔};}
\end{example}
\begin{solution}
图\ref{fig:dep-o2-match-q}中的两项在共振时反相，其模相等要求$\sqrt{R_1}=\eta\sqrt{R_2}$，即
\[
 R_1=\eta^2R_2.
\]
这就是给定输出镜与内损耗时的入射耦合匹配条件。代入透射公式，注意$q=\eta^2R_2$，有
\[
 T_{\rm res}=\frac{\eta(1-R_2)}{1-\eta^2R_2},\qquad
 A_{\rm res}=1-T_{\rm res}
 =\frac{(1-\eta)(1+\eta R_2)}{1-\eta^2R_2}.
\]
这里镜无吸收，$A_{\rm res}$全是介质内部损耗。$\eta=1$时匹配变为$R_1=R_2$，此时全部透射；$R_2=1$且$\eta<1$时没有透射口，共振入射能量全部在腔内耗散。
\ExpansionDiagram{fig:dep-o2-match-a}{共振匹配时，外部能量分成透射和内部损耗}{
\draw[->,thick](0,0)--(2,0)node[midway,above]{$P$};\draw(2,-.6)rectangle(4,.6)node[midway]{腔};\draw[->,thick](4,0)--(6,0)node[midway,above]{$TP$};\draw[->,thick](3,-.6)--(3,-1.6)node[below]{$AP$};\node[above]at(3,.8){反射为零};}
例如$\eta=0.9,R_2=0.8$，所需$R_1=0.648$，得到$T_{\rm res}=0.18/0.352\simeq0.5114$、$A_{\rm res}\simeq0.4886$。图\ref{fig:dep-o2-match-a}说明反射消失是相干抵消条件，能量的最终去向还须由透射和吸收分别计算。若镜面也吸收，应重新区分镜的透射率与$1-R_j$，不能直接使用本题的分子。
\end{solution}

\subsection{从四次强度测量还原复干涉项}
只测一次强度不能同时确定振幅和相位。给一臂施加已知附加相位$\alpha$，强度写成
\[
 I(\alpha)=B+C\cos(\phi+\alpha),\qquad B=I_1+I_2,
 \quad C=2\sqrt{I_1I_2}|\gamma|.
\]
$\gamma$是两路的复相干度，其辐角已并入$\phi$。在四个等间隔相位$0,\pi/2,\pi,3\pi/2$测量，余弦项和正弦项可以分别提取；这与把两幅条纹的位置目测相减不同。

\begin{example}四步移相与相位的象限
某点四次强度依次为$14,6,2,10$。已知附加相位顺序为$0,\pi/2,\pi,3\pi/2$。求$B,C,\phi$与可见度；若单独遮挡两臂测得$I_1=I_2=4$，求$|\gamma|$。
\ExpansionDiagram{fig:dep-o2-step-q}{参考臂的已知移相使同一测量点获得四个强度样本}{
\draw(0,0)--(1,1)--(4,1)--(5,0);\draw(0,0)--(1,-1)--(4,-1)--(5,0);\draw[thick](3,-1.35)--(3,-.65);\node[below right]at(3,-1.35){$\alpha$};\fill(5,0)circle(1.4pt)node[right]{$P$};\node[above]at(2,1){物臂};\node[below]at(1,-1){参考臂};}
\end{example}
\begin{solution}
图\ref{fig:dep-o2-step-q}中两臂功率不变，只有参考臂相位变化。用相差$\pi$的样本相加，有$B=(I_0+I_\pi)/2=8$，另一对也给$(6+10)/2=8$。相减得到
\[
 X=\frac{I_0-I_\pi}{2}=C\cos\phi=6,\qquad
 Y=\frac{I_{3\pi/2}-I_{\pi/2}}2=C\sin\phi=2.
\]
于是$C=\sqrt{X^2+Y^2}=2\sqrt{10}$，$\phi=\operatorname{atan2}(2,6)\simeq18.435^\circ$，而不是仅由一个反余弦随意选分支。
\ExpansionDiagram{fig:dep-o2-step-a}{强度差给出干涉项在复平面上的两个分量}{
\draw[->](0,0)--(4,0)node[right]{$X$};\draw[->](0,0)--(0,2.4)node[above]{$Y$};\draw[->,thick](0,0)--(3,1)node[above right]{$6+2\mathrm i$};\draw[dashed](3,0)--(3,1)--(0,1);\node[below]at(3,0){$6$};\node[left]at(0,1){$2$};}
由图\ref{fig:dep-o2-step-a}的模得到可见度$\mathcal V=C/B=\sqrt{10}/4\simeq0.7906$。单臂测量给$2\sqrt{I_1I_2}=8$，所以$|\gamma|=C/8=\sqrt{10}/4$。若不知道单臂强度，仅凭四步数据只能得$B,C$，不能区分可见度降低来自功率不等还是相干度不足。两对样本和相等是该模型的一个数据相容条件；漂移、移相误差或探测非线性会破坏它。
\end{solution}

\section{谱线形状、相位扩散与色散包络}
\subsection{相位扩散对应Lorentz线形}
在$D\ll\omega_0$的窄带近似下，用$E(t)=E_0\exp[-\mathrm i\omega_0t+\mathrm i\varphi(t)]$表示幅度近乎恒定、相位缓慢扩散的场。设相位增量为均值零的高斯变量，$\langle[\varphi(t+\tau)-\varphi(t)]^2\rangle=2D|\tau|$，初相位在整周均匀分布。高斯平均给
\[
 \langle\ee^{\mathrm i[\varphi(t+\tau)-\varphi(t)]}\rangle
 =\ee^{-D|\tau|}.
\]
高斯平均可由习题\ref{ex:o2phaseerror}的特征函数推导核对：在其中取方差$s^2=2D|\tau|$、$k=1$。于是归一化场相关函数含$\ee^{-D|\tau|}\ee^{-\mathrm i\omega_0\tau}$。令$\Omega=\omega-\omega_0$，其傅里叶积分按正负时间分开：
\begin{align*}
 S(\omega)&=\frac1{2\pi}\left[
 \int_{-\infty}^0\ee^{(D+\mathrm i\Omega)\tau}\dd\tau
 +\int_0^\infty\ee^{(-D+\mathrm i\Omega)\tau}\dd\tau\right]\\
 &=\frac1{2\pi}\left[
 \frac{\ee^{(D+\mathrm i\Omega)\tau}}{D+\mathrm i\Omega}\right]_{-\infty}^0
 +\frac1{2\pi}\left[
 \frac{\ee^{(-D+\mathrm i\Omega)\tau}}{-D+\mathrm i\Omega}\right]_0^\infty\\
 &=\frac1{2\pi}\left(\frac1{D+\mathrm i\Omega}
 +\frac1{D-\mathrm i\Omega}\right).
\end{align*}
通分后得到归一化角频率谱
\[
 S(\omega)=\frac{D}{\pi[(\omega-\omega_0)^2+D^2]}.
\]
角频率半高全宽为$2D$，普通频率半高全宽为$\Delta\nu=D/\pi$。等强两臂的可见度为$\ee^{-D|\tau|}$，不存在矩形谱那样的有限首零点。若将相干长度定义为可见度降至$1/\ee$的光程差，则$L_c=c/D=c/(\pi\Delta\nu)$。定义换为半高宽时常数会改变，应连同定义一起报告。

\begin{example}光源线宽与扫描曝光共同降低可见度
等强双臂干涉仪采用上述相位扩散光源，中心波长$600\,\mathrm{nm}$，谱线全宽$\Delta\nu=3.0\,\mathrm{MHz}$。两臂光程差$10\,\mathrm m$。扫描镜以$30\,\mu\mathrm{m/s}$匀速运动，使往返光程差以$2v$改变，曝光时间为$5.0\,\mathrm{ms}$。在曝光期间相关包络近似不变，求记录的条纹可见度。
\ExpansionDiagram{fig:dep-o2-diffusion-q}{扫描镜改变载波相位，光源线宽另行决定延迟相关包络}{
\draw[->](-1.8,0)--(-.25,0);\draw[thick](-.2,-.2)--(.2,.2);\draw[->](0,.2)--(0,2);\draw[->](.12,2)--(.12,.2);\draw[very thick](-.4,2)--(.5,2)node[right]{固定镜};\draw[->](.25,0)--(4,0);\draw[->](4,-.12)--(.2,-.12);\draw[very thick](4,-.6)--(4,.6);\draw[->](4,.9)--(5.2,.9)node[right]{$v$};\draw[->](0,-.3)--(0,-1.5)node[below]{记录点};}
\end{example}
\begin{solution}
图\ref{fig:dep-o2-diffusion-q}中静止光程差给$\tau=10/c$，所以包络因子为$\exp[-\pi\Delta\nu\,10/c]=\exp(-\pi/10)\simeq0.7304$，取$c=3.0\times10^8\,\mathrm{m/s}$。

镜面移动使相位角速度$\Omega_s=4\pi v/\lambda$。在以曝光中点为原点的时间区间内平均余弦，
\[
 \frac1{T_e}\int_{-T_e/2}^{T_e/2}\cos(\phi_0+\Omega_st)\dd t
 =\left[\frac{\sin(\phi_0+\Omega_st)}{\Omega_sT_e}\right]_{-T_e/2}^{T_e/2}
 =\operatorname{sinc}(\Omega_sT_e/2)\cos\phi_0.
\]
本题$\Omega_sT_e/2=2\pi vT_e/\lambda=\pi/2$，所以记录可见度为
\[
 \mathcal V=\ee^{-\pi/10}\frac2\pi\simeq0.4650.
\]
\ExpansionDiagram{fig:dep-o2-diffusion-a}{可见度的两个独立因子：延迟包络与曝光平均}{
\node[draw]at(0,0){$0.7304$};\node at(1.5,0){$\times$};\node[draw]at(3,0){$2/\pi$};\node at(4.5,0){$=$};\node[draw]at(6,0){$0.4650$};\node[below]at(0,-.4){谱线宽度};\node[below]at(3,-.4){扫描曝光};}
图\ref{fig:dep-o2-diffusion-a}中的两项分别来自不同平均。本题曝光内包络变化参数$D(2vT_e/c)$远小于1，可以近似相乘；若这个量不可略，应把随时间改变的相关包络一并放入曝光积分。减小曝光时间只能改善后一项，不能恢复已经因长光程差而失去的时间相干性。
\end{solution}

\subsection{群光程决定包络，中心相位决定细条纹}
白光或短脉冲干涉中，两臂相位差通常依赖频率。设中心频率附近$\Omega=\omega-\omega_0$，相位展开为
\[
 \phi(\omega)=\phi_0+\tau_g\Omega+\tfrac12D_2\Omega^2,
\]
入射功率谱的归一化形状为$\exp[-\Omega^2/(2\sigma_\omega^2)]/(\sqrt{2\pi}\sigma_\omega)$。取$\sigma_\omega\ll\omega_0$，将偏移频率积分延伸到整个实轴。复干涉项是$\langle\ee^{\mathrm i\phi(\omega)}\rangle$。记$a=(1-\mathrm iD_2\sigma_\omega^2)/(2\sigma_\omega^2)$、$b=\mathrm i\tau_g$，则$\operatorname{Re}a>0$，并有
\[
 \Gamma=\frac{\ee^{\mathrm i\phi_0}}{\sqrt{2\pi}\sigma_\omega}
 \int_{-\infty}^{\infty}\ee^{-a\Omega^2+b\Omega}\dd\Omega
 =\frac{\ee^{\mathrm i\phi_0}}{\sqrt{2\pi}\sigma_\omega}
 \sqrt{\frac\pi a}\ee^{b^2/(4a)}.
\]
最后一步使用式\eqref{eq:o1gaussintegral}已推导的积分。代入$a,b$，得
\[
 \Gamma=\frac{\ee^{\mathrm i\phi_0}}{\sqrt{1-\mathrm iD_2\sigma_\omega^2}}
 \exp\!\left[-\frac{\sigma_\omega^2\tau_g^2}
 {2(1-\mathrm iD_2\sigma_\omega^2)}\right].
\]
平方根取在$D_2=0$处连续为1的分支。模为
\[
 |\Gamma|=[1+(D_2\sigma_\omega^2)^2]^{-1/4}
 \exp\!\left[-\frac{\sigma_\omega^2\tau_g^2}
 {2[1+(D_2\sigma_\omega^2)^2]}\right].
\]
二阶色散使包络展宽、峰值降低。包络最大值位于$\tau_g=0$，细条纹相位却还含有$\phi_0$和积分产生的附加相位，因此包络峰的位置一般不等于中心频率的零程差位置。

\begin{example}插入色散片后应补偿哪一种光程
单程双臂干涉仪在一臂垂直插入厚度$l$、折射率$n(\omega)$的透明片，另一臂增加空气光程$x$。忽略反射，近中心频率相位差为$\phi(\omega)=\omega[l(n(\omega)-1)-x]/c$。求使相关包络最大的一阶补偿量，与使中心相位为零的补偿量比较。
\ExpansionDiagram{fig:dep-o2-group-q}{一臂色散片、另一臂空气光程补偿}{
\draw(0,0)--(1,1.2)--(4,1.2)--(5,0);\draw(0,0)--(1,-1.2)--(4,-1.2)--(5,0);\draw[fill=gray!15](2.2,.9)rectangle(2.6,1.5);\node[above]at(2.4,1.6){$l,n(\omega)$};\node[below]at(2.5,-1.3){空气补偿$x$};}
\end{example}
\begin{solution}
图\ref{fig:dep-o2-group-q}的中心相位为$\phi_0=\omega_0[l(n_0-1)-x]/c$。对完整的$\omega n(\omega)$求导，
\[
 \tau_g=\left.\frac{\dd\phi}{\dd\omega}\right|_{\omega_0}
 =\left.\frac{l[n(\omega)-1]-x+\omega l n'(\omega)}c\right|_{\omega_0}
 =\frac{l[n_0+\omega_0n'_0-1]-x}{c}.
\]
因此包络最大要求$x_g=l(n_g-1)$，其中$n_g=n_0+\omega_0n'_0$；中心频率相位恰为零则要求$x_p=l(n_0-1)$。两者差为$x_g-x_p=l\omega_0n'_0$。
\ExpansionDiagram{fig:dep-o2-group-a}{细条纹位于相关包络之内，零载波相位不必位于包络中心}{
\draw[->](0,0)--(7,0)node[right]{$x$};\draw[thick,domain=.1:6.6,samples=130]plot(\x,{1.2*exp(-.35*(\x-3.8)^2)});\draw[dashed,domain=.1:6.6,samples=200]plot(\x,{1.2*exp(-.35*(\x-3.8)^2)*cos(600*(\x-2.5))});\draw[dotted](2.5,-1.2)--(2.5,1.6)node[above]{$x_p$};\draw[dotted](3.8,-1.2)--(3.8,1.6)node[above]{$x_g$};}
图\ref{fig:dep-o2-group-a}示意两种中心的位置，细条纹的精确位置还由谱积分相位决定。忽略色散$n'_0=0$时，两种补偿相同。二阶色散显著时，改变$x$只调整常数项和一次频率项，$D_2$仍然保留；恢复原有峰值需要另一种材料或其他色散补偿，增加采样密度对这一展宽没有作用。
\end{solution}

\section{习题与解答}
\begin{exercise}\label{ex:o2bragg}
正常入射的无损多层反射镜由$N$组高、低折射率层$H,L$组成，入射侧先遇$H$，每层在设计波长的光学厚度均为$\lambda_0/4$。入射介质、基底导纳分别为$n_0,n_s$。推导设计波长处的反射率。
\Diagram{fig:o2braggq}{交替四分之一波厚层叠加界面反射}{\foreach\x in{0,1,2,3}{\draw[curve](\x,-1)rectangle({\x+.6},1);\node[above]at({\x+.3},1){$H$};\draw[guide]({\x+.6},-1)rectangle({\x+1},1);\node[below]at({\x+.8},-1){$L$};}\draw[axis](-1,0)--(5,0);}
\end{exercise}
\begin{solution}
把磁场归一化为$\mathcal H=Z_0H$，其中$Z_0$为真空波阻抗。正常入射层中前、后行波幅为$A,B$，在某参考面有$E=A+B$、$\mathcal H=n(A-B)$。厚度为$l$的层内，前行波向右乘$\ee^{i\phi}$、后行波向右乘$\ee^{-i\phi}$，$\phi=k_0nl$。因此左面振幅可由右面振幅写为$A_L=A_R\ee^{-i\phi}$、$B_L=B_R\ee^{i\phi}$。代入场的和、差，得到从层右面到左面的特征矩阵
\[
 \begin{pmatrix}E_L\\\mathcal H_L\end{pmatrix}
 =C(n,\phi)\begin{pmatrix}E_R\\\mathcal H_R\end{pmatrix},\qquad
 C(n,\phi)=\begin{pmatrix}\cos\phi&-i\sin\phi/n\\
 -in\sin\phi&\cos\phi\end{pmatrix}.
\]
界面处切向$E,\mathcal H$连续，可以直接接续相邻层。图\ref{fig:o2braggq}在设计波长$\phi_H=\phi_L=\pi/2$，一组矩阵
\[
 C_HC_L=\begin{pmatrix}-n_L/n_H&0\\0&-n_H/n_L\end{pmatrix}.
\]
设$q=n_H/n_L$，$N$组矩阵为$(-1)^N\operatorname{diag}(q^{-N},q^N)$。基底中只有向右的波，故$\mathcal H_s=n_sE_s$，于是$E_{\rm in}=(-1)^Nq^{-N}E_s$、$\mathcal H_{\rm in}=(-1)^Nq^Nn_sE_s$，等效输入导纳为$Y=\mathcal H_{\rm in}/E_{\rm in}=n_sq^{2N}$。
\Diagram{fig:o2bragga}{层组把基底导纳变换为越来越大的输入导纳}{\draw[curve](0,0)rectangle(2.8,1);\node at(1.4,.5){$n_s$};\draw[->](3,.5)--(4,.5);\node[right]at(4.2,.5){$Y=n_s(n_H/n_L)^{2N}$};\node[below]at(2.3,-.7){$r=(n_0-Y)/(n_0+Y)$};}
入射和反射叠加给$E=1+r$、$\mathcal H=n_0(1-r)$，与$\mathcal H=YE$联立得到图\ref{fig:o2bragga}的$r$，反射率$R=[(n_0-Y)/(n_0+Y)]^2$。当$n_H>n_L$且$N$增大，$R\to1$。该结果只在设计波长和给定层序精确成立；交换层序会倒置导纳倍率，层数和层序均应保留。
\end{solution}

\begin{exercise}\label{ex:o2jointcoherence}
均匀线源宽$w$、距双缝$a$，双缝间距$d$，光谱为中心$\nu_0$、宽$\Delta\nu$的矩形谱。源位置与频率统计独立，两路等强、偏振相同，中心源点的两路时间差为$\tau$。在窄带近似下求可见度，并说明何时不能简单把时间与空间可见度相乘。
\Diagram{fig:o2jointq}{源位置和频率各自造成一类相位散布}{\draw[curve](0,0)rectangle(2.2,1.2);\node at(1.1,.6){位置$x_s$};\draw[curve](3.6,0)rectangle(5.8,1.2);\node at(4.7,.6){频率$\nu$};\node[below]at(2.9,-.7){两变量分布独立};}
\end{exercise}
\begin{solution}
图\ref{fig:o2jointq}中各源元、各频率的两路相位差为$2\pi\nu\tau+2\pi\nu d x_s/(ca)$。对均匀位置分布和矩形谱，归一化干涉因子为
\[
 \Gamma=\frac1{w\Delta\nu}\int_{-w/2}^{w/2}\dd x_s
 \int_{\nu_0-\Delta\nu/2}^{\nu_0+\Delta\nu/2}
 \ee^{i2\pi\nu\tau+i2\pi\nu d x_s/(ca)}\dd\nu.
\]
窄带时在位置项中以$\nu_0$代替$\nu$，积分核分解。频率积分给$\ee^{i2\pi\nu_0\tau}\operatorname{sinc}(\pi\Delta\nu\tau)$，位置积分给$\operatorname{sinc}[\pi\nu_0dw/(ca)]$。假定两路等强且偏振相同，可见度为$|\Gamma|$，因此
\[
 \mathcal V\simeq\left|\operatorname{sinc}(\pi\Delta\nu\tau)
 \operatorname{sinc}\left(\frac{\pi dw}{\lambda_0a}\right)\right|.
\]
频率平均控制不同延迟的同点相干，位置平均控制两缝的照明相干，两者都要满足才有高对比条纹。
\Diagram{fig:o2jointa}{两个独立平均因子相乘，任一个接近零都会洗去条纹}{\node at(0,.6){时间因子};\node at(0,-.6){空间因子};\draw[->](1.4,.6)--(3,0);\draw[->](1.4,-.6)--(3,0);\node[right]at(3.2,0){总可见度};}
被略去的交叉相位尺度为$2\pi\Delta\nu d w/(ca)$，须远小于1。若谱很宽、源很大，或不同频率来自不同源位置，图\ref{fig:o2jointa}的分解不成立，应保留联合分布做二重积分。两个单独可见度接近1只是窄带独立模型中的充分描述，不能对任意源结构机械相乘。
\end{solution}

\begin{exercise}\label{ex:o2coherencenum}
中心波长600 nm、窄矩形谱宽度0.2 nm的光源用于双束干涉。估计首零光程差及迈克耳孙单镜相对零光程差允许的首零移动量。
\Diagram{fig:o2cohnumq}{镜移通过往返光程差扫描相干包络}{\draw[very thick](1,-1)--(1,1);\draw[guide](3,-1)--(3,1);\draw[<->](1,-1.5)--(3,-1.5)node[midway,below]{$d$};\draw[->](0,.4)--(2.7,.4);\draw[->](2.7,-.3)--(0,-.3);}
\end{exercise}
\begin{solution}
由$\nu=c/\lambda$求导，$\dd\nu=-c\dd\lambda/\lambda^2$；谱宽取正值，窄谱时$\Delta\nu\simeq c\Delta\lambda/\lambda_0^2$。图\ref{fig:o2cohnumq}的首零光程差为$\mathcal D_0=c/\Delta\nu\simeq\lambda_0^2/\Delta\lambda=(600\,\mathrm{nm})^2/(.2\,\mathrm{nm})=1.8$ mm。
\Diagram{fig:o2cohnuma}{单镜移动量是扫描光程差的一半}{\draw[axis](0,0)--(5,0)node[right]{毫米};\draw[curve](1,-.4)--(1,.4);\draw[curve](2.8,-.4)--(2.8,.4);\node[below]at(1,-.4){$d=.9$};\node[below]at(2.8,-.4){$\mathcal D=1.8$};}
图\ref{fig:o2cohnuma}单镜首零移动量$d_0=.9$ mm。此数以光程差零点为起点；若初始已有臂长差，应把它计入，只有从零光程差开始移动，才在.9 mm处首次消失。矩形谱理想模型在首零以外仍有较小旁瓣，可见度不是永久为零。
\end{solution}

\begin{exercise}\label{ex:o2sourceorientation}
均匀线光源长度$w$与双缝连线成角$\psi$，距缝$a$。求空间相干可见度，比较光源沿缝长方向和垂直缝长方向的情况。
\Diagram{fig:o2orientationq}{光源的基线投影决定两缝之间的入射相位散布}{\draw[axis](-.3,0)--(3.7,0)node[right]{双缝基线};\draw[very thick](0,0)--(2.2,1.7);\draw[guide](2.2,0)--(2.2,1.7);\node at(.8,.3){$\psi$};\node[right]at(2.2,1.7){$w$};}
\end{exercise}
\begin{solution}
图\ref{fig:o2orientationq}线源参数$s\in[-w/2,w/2]$，其在双缝基线上的投影为$s\cos\psi$，源元相位为$2\pi d s\cos\psi/(\lambda a)$。将式\eqref{eq:o2rectaverage}中的相位斜率改为$2\pi d\cos\psi/(\lambda a)$，区间长度仍取$w$，得到$\mathcal V=|\operatorname{sinc}(\pi dw\cos\psi/(\lambda a))|$；$\cos\psi=0$时用该定积分的常相位分支，结果为1。
\Diagram{fig:o2orientationa}{垂直基线延伸时，近轴模型中两缝相位差不随源点改变}{\fill(-1,0)circle(1.5pt)node[below]{$S_-$};\fill(1,0)circle(1.5pt)node[below]{$S_+$};\draw[very thick](0,.7)--(0,2.5);\draw[guide](0,2)--(-1,0);\draw[guide](0,2)--(1,0);\node[right]at(.2,2.5){源的投影示意};}
沿缝长方向通常垂直双缝基线，$\cos\psi=0$，图\ref{fig:o2orientationa}的源点不引起一阶基线相位差，可见度保持1。沿基线延伸最容易洗掉条纹。实际有限孔径和非近轴高阶项仍可限制可见度，不能把此一阶结论外推为任意大光源都完全相干。
\end{solution}

\begin{exercise}\label{ex:o2phaseerror}
两束等强光相位差含随机误差$\epsilon$，服从均值0、方差$s^2$的高斯分布。求长曝光可见度，解释相位抖动对条纹位置和对比度的不同影响。
\Diagram{fig:o2errorq}{相位随时间在平均位置附近抖动}{\draw[axis](0,0)--(5,0)node[right]{$t$};\draw[axis](0,-1.4)--(0,1.7)node[above]{$\epsilon$};\draw[curve](0,.2)--(.5,-.6)--(1,.8)--(1.5,.2)--(2,-.9)--(2.5,.6)--(3,-.2)--(3.5,1)--(4,.1)--(4.5,-.4);}
\end{exercise}
\begin{solution}
图\ref{fig:o2errorq}中，瞬时强度可写为$2I_1[1+\cos(\delta_0+\epsilon)]$。长曝光的交叉项由$G=\langle\ee^{i\epsilon}\rangle$控制，其中
\[
 G=\frac1{\sqrt{2\pi}s}\int_{-\infty}^\infty
 \exp[-\epsilon^2/(2s^2)]\ee^{i\epsilon}\dd\epsilon.
\]
引入$G(k)=\langle\ee^{ik\epsilon}\rangle$。概率密度的导数为$p'(\epsilon)=-(\epsilon/s^2)p(\epsilon)$；对$k$求导，再分部积分，
\begin{align*}
 G'(k)&=i\int_{-\infty}^{\infty}\epsilon p(\epsilon)\ee^{ik\epsilon}\dd\epsilon
 =-is^2\int_{-\infty}^{\infty}p'(\epsilon)\ee^{ik\epsilon}\dd\epsilon\\
 &=-is^2[p(\epsilon)\ee^{ik\epsilon}]_{-\infty}^{\infty}
 +is^2\int_{-\infty}^{\infty}ikp(\epsilon)\ee^{ik\epsilon}\dd\epsilon
 =-s^2kG(k).
\end{align*}
边界项因高斯衰减而为零。以$G(0)=1$分离变量积分，$\int_1^{G(k)}\dd g/g=-s^2\int_0^kq\dd q=-s^2k^2/2$，因此$G(k)=\ee^{-s^2k^2/2}$。
\Diagram{fig:o2errora}{零均值抖动不移动平均条纹，却降低可见度}{\draw[axis](0,0)--(4.5,0)node[right]{$s$};\draw[axis](0,0)--(0,2.7)node[above]{$\mathcal V$};\draw[curve,domain=0:4,samples=60]plot(\x,{2.3*exp(-\x*\x/2)});}
取$k=1$得到图\ref{fig:o2errora}的平均强度$2I_1[1+\ee^{-s^2/2}\cos\delta_0]$，因此$\mathcal V=\ee^{-s^2/2}$。平均相位误差为零时亮暗位置不变；若均值为$\bar\epsilon$，则$G=\ee^{i\bar\epsilon}\ee^{-s^2/2}$，条纹另有整体位移。

这里$s$以弧度计，曝光须覆盖足够多次抖动，且误差统计与光强无关。若机械位移$z$使迈克耳孙一臂镜往返程差改变$2z$，则相位标准差为$s=4\pi\sigma_z/\lambda$。同一位置抖动对短波长的对比度损失更大，不能只用$\sigma_z$而不换成相位来比较。
\end{solution}

\begin{exercise}\label{ex:o2filmangle}
一薄膜反射具有一次相对半波损失。正常入射某波长$\lambda_0$在膜厚$h$处形成$m$级暗纹。入射角改为$i$，求对应同级暗纹波长，介质色散忽略。
\Diagram{fig:o2angleq}{斜入射改变膜内法向相位积累}{\draw[curve](-2,0)--(2,0);\draw[curve](-2,-.7)--(2,-.7);\draw[->](-1.6,1.6)--(0,0)--(.4,-.7)--(.8,0)--(2.4,1.6);\node[left]at(-2,-.35){$n$};}
\end{exercise}
\begin{solution}
图\ref{fig:o2angleq}正常暗纹条件为$2nh=m\lambda_0$。对正膜厚$h>0$，取$m\geq1$；斜入射时$2nh\cos r=m\lambda_i$，且$\sin r=\sin i/n$，两暗纹条件相除，得$\lambda_i=\lambda_0\sqrt{1-\sin^2i/n^2}$。若$h=0,m=0$，接触处在该反射相位模型下对所有波长都暗，不能由零级条件选出唯一的$\lambda_i$。
\Diagram{fig:o2anglea}{同级暗纹波长随角度增加向短波移动}{\draw[axis](0,0)--(4.4,0)node[right]{$i$};\draw[axis](0,0)--(0,2.7)node[above]{$\lambda_i/\lambda_0$};\draw[curve,domain=0:85,samples=60]plot({\x/22},{2.3*sqrt(1-sin(\x)^2/2.25)});\node[left]at(0,2.3){$1$};}
图\ref{fig:o2anglea}对应$n=1.5$。角度增大使法向波矢减小，暗纹波长蓝移。若两界面反射相位在所扫角度内因$p$偏振过布儒斯特点而改变，原来的暗纹编号规则需要重新讨论，不能只连续追踪$\cos r$。
\end{solution}

\begin{exercise}\label{ex:o2freespectral}
平行腔长度$l$、折射率$n$无色散，正常入射。求相邻透射共振的频率间隔，并由精细度$\mathcal F_c$求线宽。
\Diagram{fig:o2fsrq}{共振要求往返相位为$2m\pi$}{\draw[very thick](0,-1)--(0,1);\draw[very thick](4,-1)--(4,1);\draw[->](.2,.4)--(3.8,.4);\draw[->](3.8,-.4)--(.2,-.4);\node at(2,0){$n,l$};}
\end{exercise}
\begin{solution}
图\ref{fig:o2fsrq}的单程相位为$\phi=2\pi nl\nu/c$，共振要求$2\phi=2\pi m$，因此
\[
 \nu_m=\frac{mc}{2nl},\qquad
 \Delta\nu_{\rm FSR}=\nu_{m+1}-\nu_m=\frac c{2nl}.
\]
精细度按相邻峰距除以单峰半高全宽定义，故$\delta\nu=\Delta\nu_{\rm FSR}/\mathcal F_c$。高反射率时再代入$\mathcal F_c\simeq\pi\sqrt R/(1-R)$。
\Diagram{fig:o2fsra}{峰间距与峰宽是不同的两个频率尺度}{\draw[axis](0,0)--(5.5,0)node[right]{$\nu$};\foreach\x in{1,3,5}{\draw[curve]({\x-.3},0)..controls({\x-.1},2.8)and({\x+.1},2.8)..({\x+.3},0);}\draw[<->](1,-.5)--(3,-.5)node[midway,below]{FSR};}
图\ref{fig:o2fsra}中增大$R$使单峰变窄，而相邻峰距仍由光学长度决定。若$n$有色散，往返相位应写$\Psi(\nu)=4\pi l\nu n(\nu)/c$，近邻共振满足
\[
 2\pi\simeq\Psi'(\nu)\Delta\nu,
 \quad \Delta\nu\simeq\frac{c}{2l[n+\nu\dd n/\dd\nu]}
 =\frac c{2ln_g}.
\]
其中$n_g=n-\lambda\dd n/\dd\lambda$是群折射率。若镜的反射相位也随频率变化，还须把两个镜面的反射相位导数加进$\Psi'$；此时有效往返群延迟不再只是$2nl/c$。只要相邻峰之间参数变化足够小，上述导数法仍适用。
\end{solution}

\begin{exercise}\label{ex:o2Nphasechange}
$N$束等振幅光相邻相位差由$\delta$均匀增到$\delta+\epsilon$。求主峰中心的相位移动与峰高变化；若只在最后一束附加小相位$\epsilon$，轴上峰高的二阶变化如何？
\Diagram{fig:o2Nphaseq}{线性相位斜坡与单束相位误差有不同后果}{\draw[axis](0,0)--(4.5,0)node[right]{$j$};\draw[axis](0,0)--(0,2.7)node[above]{附加相位};\draw[curve](0,0)--(3.8,2);\draw[guide](0,0)--(3.8,0)--(3.8,1);}
\end{exercise}
\begin{solution}
图\ref{fig:o2Nphaseq}线性斜坡只把多束强度函数自变量换成$\delta+\epsilon$，主峰移到$\delta=2m\pi-\epsilon$，峰高仍$N^2I_1$。若只有最后一束出错，原轴上合场为$A[(N-1)+\ee^{i\epsilon}]$。
\Diagram{fig:o2Nphasea}{单束误差使一条相量偏离其余同向相量}{\draw[->,very thick](0,0)--(3.3,0)node[below]{$(N-1)A$};\draw[->,thick](3.3,0)--(4.1,1)node[right]{$A\ee^{i\epsilon}$};\draw[guide](3.3,0)--(4.5,0);}
图\ref{fig:o2Nphasea}给$I/I_1=(N-1)^2+1+2(N-1)\cos\epsilon\simeq N^2-(N-1)\epsilon^2$。这是原轴上强度，严格新峰位置还需对$\delta$求极值；不能把它直接当全部角度中已重新寻优的最大值。相位斜坡可用峰位补偿，局部误差一般同时改变旁瓣和峰形。
\end{solution}

\begin{exercise}\label{ex:o2visibility}
两相干光强比为9:1，偏振方向夹角$\eta$，求条纹可见度。若希望完全相消，需同时满足什么条件？
\Diagram{fig:o2visibilityq}{两束电场的偏振方向夹角}{\draw[axis](0,0)--(3,0)node[right]{$\vct e_1$};\draw[axis](0,0)--(1.5,2.2)node[right]{$\vct e_2$};\node at(.8,.4){$\eta$};}
\end{exercise}
\begin{solution}
以图\ref{fig:o2visibilityq}的第一束偏振为$x$轴，令两束实振幅为$3A,A$。第二束可分成$A\cos\eta$的平行分量和$A\sin\eta$的正交分量。合场分量为
\[
 E_x=3A+A\cos\eta\ee^{i\delta},\qquad
 E_y=A\sin\eta\ee^{i\delta}.
\]
探测器记录两正交分量强度之和，因此$I/I_A=|E_x/A|^2+|E_y/A|^2=10+6\cos\eta\cos\delta$。
\Diagram{fig:o2visibilitya}{正交分量不能被另一束光抵消}{\draw[axis](0,0)--(2.4,0)node[right]{$\parallel\vct e_1$};\draw[axis](0,0)--(0,2.2)node[above]{垂直分量};\draw[->,thick](0,0)--(1.4,1.6);\draw[guide](1.4,0)--(1.4,1.6)--(0,1.6);}
图\ref{fig:o2visibilitya}中的正交分量不能与第一束相消。相位变化时，$I_{\max}/I_A=10+6|\cos\eta|$、$I_{\min}/I_A=10-6|\cos\eta|$，所以$\mathcal V=3|\cos\eta|/5$。即使偏振平行，最暗仍为$4I_A$，因为$3A-A=2A$。

一般两束完全相干线偏振光的可见度为$2\sqrt{I_1I_2}|\cos\eta|/(I_1+I_2)$。由$2\sqrt{I_1I_2}\leq I_1+I_2$与$|\cos\eta|\leq1$，完全暗点要求两不等式同时取等号：两束等强、偏振共线，再选择相反的复场相位。若把偏振矢量方向反转，振幅相位也随之改变$\pi$，应以两束总电场反向为最终判据。部分相干时交叉项还要乘$|\gamma_{12}|$。
\end{solution}

\begin{exercise}\label{ex:o2filmindex}
空气—薄膜—基底的折射率依次为$1,n_f,n_s$，$1<n_f<n_s$。垂直照射，求使反射最弱的最小膜厚；进一步求两反射振幅可完全抵消的折射率条件。
\Diagram{fig:o2antirefq}{法向入射时两反射同向；分列画出各束以显示传播次序}{
\draw[curve](-2,0)--(2.4,0);\draw[curve](-2,-.9)--(2.4,-.9);
\node[left]at(-2,1.25){$1$};\node[left]at(-2,-.45){$n_f$};\node[left]at(-2,-1.3){$n_s$};
\draw[->](-1.2,1.5)--(-1.2,0)node[midway,left]{入射};
\draw[->](-.65,0)--(-.65,1.5)node[above]{$r_{01}$};
\draw[->](.35,.2)--(.35,-.9);
\draw[->](.9,-.9)--(.9,1.5)node[above]{$r_{12}\ee^{2i\phi}$};
\draw[<->](1.9,0)--(1.9,-.9)node[midway,right]{$h$};
}
\end{exercise}
\begin{solution}
图\ref{fig:o2antirefq}中两界面均由低折射率向高折射率反射，$r_{01},r_{12}$都为负；两反射的相对初相位为零。要使两束因传播反相，最小正厚度满足$2n_fh=\lambda_0/2$，故$h=\lambda_0/(4n_f)$。这给出相位条件，完全消反还需振幅匹配。

令单程膜内相位为$\phi$。纳入多次反射，首个由下界面返回的振幅为$t_{01}t_{10}r_{12}\ee^{2i\phi}$，后续倍率为$r_{10}r_{12}\ee^{2i\phi}$。因$r_{10}=-r_{01}$、$t_{01}t_{10}=1-r_{01}^2$，
\[
 r=r_{01}+\frac{(1-r_{01}^2)r_{12}\ee^{2i\phi}}
 {1+r_{01}r_{12}\ee^{2i\phi}}
 =\frac{r_{01}+r_{12}\ee^{2i\phi}}{1+r_{01}r_{12}\ee^{2i\phi}}.
\]
\Diagram{fig:o2antirefa}{四分之一波膜让往返传播相位成为$\pi$}{\draw[axis](0,0)--(2.7,0)node[right]{$r_{01}$};\draw[axis](0,0)--(-2.7,0)node[left]{$r_{12}\ee^{i\pi}$};\node[above]at(0,1){两反射相量等大反向时消失};}
在图\ref{fig:o2antirefa}的四分之一波厚度，$\ee^{2i\phi}=-1$，消反要求$r_{01}=r_{12}$。代入界面系数并交叉相乘，
\[
 (1-n_f)(n_f+n_s)=(n_f-n_s)(1+n_f)
 \quad\Longrightarrow\quad n_f^2=n_s.
\]
因此$n_f=\sqrt{n_s}$。若入射介质改为$n_0$，相同推导给$n_f^2=n_0n_s$。匹配只针对给定波长和法向入射；色散或斜入射会改变相位与两偏振的导纳。若折射率不匹配，四分之一波厚度仍使反射最小，但最低反射率为$[(r_{01}-r_{12})/(1-r_{01}r_{12})]^2$，一般不为零。
\end{solution}

\begin{exercise}\label{ex:o2ringliquid}
同一牛顿环装置在空气和未知液体中用同一波长测得相隔$p$级暗环的直径平方差分别为$D_a,D_l$。求液体折射率，说明如何消去中心小间隙影响。
\Diagram{fig:o2liquidq}{同一球面间隙先后由空气和液体填充}{\draw[curve](-2.3,0)--(2.3,0);\draw[curve,domain=-2:2,samples=50]plot(\x,{.15+.18*\x*\x});\node at(0,.6){$n$};\node[above]at(0,1.4){几何$R$不变};}
\end{exercise}
\begin{solution}
图\ref{fig:o2liquidq}中暗环直径$d_m=2r_m$，所以$d_{m+p}^2-d_m^2=4pR\lambda_0/n$。空气近似$n_a=1$，两次相除得$n=D_a/D_l$。
\Diagram{fig:o2liquida}{直径平方差消除所有环共同的厚度偏移}{\draw[axis](0,0)--(4.5,0)node[right]{级次$m$};\draw[axis](0,-.5)--(0,3)node[above]{$d_m^2$};\draw[curve](.7,0)--(3.8,2.6);\draw[curve](.7,0)--(4.2,1.6);\node[right]at(3.8,2.6){空气};\node[right]at(4.2,1.6){液体};}
中心间隙$h_0$只使$d_m^2$增加共同截距$-8Rh_0$，图\ref{fig:o2liquida}的斜率不受影响。如果充液引起透镜受力、曲率或接触状态改变到不能视为同一$R$，该消元就不成立；应使用同一机械装置、相同测量级差。
\end{solution}

\begin{exercise}\label{ex:o2doubleline}
迈克耳孙干涉仪受两条等强、互不相干的近邻谱线$\lambda$和$\lambda+\Delta\lambda$照射。求随光程差变化的条纹可见度包络，以及相邻可见度消失位置之间的镜移量。
\Diagram{fig:o2doubleq}{两频率各形成条纹，记录器相加强度}{\draw[curve](0,0)rectangle(2,1);\node at(1,.5){两条谱线};\draw[->](2,.5)--(3,.5);\draw[curve](3,0)rectangle(5.4,1);\node at(4.2,.5){光程差$\mathcal D$};}
\end{exercise}
\begin{solution}
取每条谱线的条纹强度为$I_0[1+\cos(k_i\mathcal D)]$。图\ref{fig:o2doubleq}的和为$2I_0[1+\cos((k_1-k_2)\mathcal D/2)\cos((k_1+k_2)\mathcal D/2)]$，所以包络可见度为$|\cos(\Delta k\mathcal D/2)|$。
\Diagram{fig:o2doublea}{快速干涉条纹受慢变化拍频包络调制}{\draw[axis](0,0)--(6.5,0)node[right]{$\mathcal D$};\draw[curve,domain=0:6,samples=170]plot(\x,{1+cos(60*\x)*cos(720*\x)});\draw[guide,domain=0:6,samples=80]plot(\x,{1+abs(cos(60*\x))});\draw[guide,domain=0:6,samples=80]plot(\x,{1-abs(cos(60*\x))});}
图\ref{fig:o2doublea}相邻包络零点相差$\Delta\mathcal D=2\pi/|\Delta k|$。由于$|\Delta k|\simeq2\pi|\Delta\lambda|/\lambda^2$，得到$\Delta\mathcal D\simeq\lambda^2/|\Delta\lambda|$。单镜位移是光程差改变的一半，所以$\Delta d\simeq\lambda^2/(2|\Delta\lambda|)$。不能把两条不同频率的复电场长期相干相加来计算静态条纹。
\end{solution}

\begin{exercise}\label{ex:o2Nwidth}
有$N=100$条等间距相干缝，周期$d=10\,\mu\mathrm m$，波长$500\,\mathrm{nm}$。求一级主极大角度、首零半宽及瑞利可分辨的最小波长差。
\Diagram{fig:o2Nq}{有限数目缝的主峰角宽}{\draw[curve](0,-1.3)--(0,1.3);\foreach\y in{-1,-.5,0,.5,1}{\draw[->](0,\y)--(3,{\y+.8});}\node[above]at(0,1.5){$N$缝};\draw[guide](0,0)--(4,0);\node[right]at(3,.8){$\theta_1$};}
\end{exercise}
\begin{solution}
图\ref{fig:o2Nq}的相邻缝相位差为$\delta=2\pi d\sin\theta/\lambda$。一级主极大对应$\delta=2\pi$，所以$\sin\theta_1=.5\,\mu\mathrm m/(10\,\mu\mathrm m)=.05$、$\theta_1=2.866^\circ$。在主峰附近令$\delta=2\pi+\epsilon$，阵列因子首零满足$N\epsilon/2=\pm\pi$，即$|\epsilon|=2\pi/N$。对角度线性展开，
\[
 \Delta\theta_0=\frac{\lambda}{Nd\cos\theta_1}
 =5.006\times10^{-4}\ \mathrm{rad}.
\]
\Diagram{fig:o2Na}{主峰到相邻首零的距离与另一谱线峰位比较}{\draw[axis](-3,0)--(3.7,0)node[right]{角度偏移};\draw[curve,domain=-2.7:-.01,samples=100]plot(\x,{2.3*(sin(180*\x)/(pi*\x))^2});\draw[curve,domain=.01:2.7,samples=100]plot(\x,{2.3*(sin(180*\x)/(pi*\x))^2});\draw[curve](-.01,2.29924)--(0,2.3)--(.01,2.29924);\draw[guide](1,0)--(1,2.3);\node[below]at(0,0){$\theta_1$};\node[below]at(1,0){首零};}
图\ref{fig:o2Na}中的瑞利条件要求一条谱线的主峰落在另一条首零处。同级角色散由$d\sin\theta=m\lambda$给$\Delta\theta\simeq m\Delta\lambda/(d\cos\theta)$。令它等于首零半宽，得到
\[
 \frac{m\Delta\lambda}{d\cos\theta}
 =\frac{\lambda}{Nd\cos\theta},\qquad
 \Delta\lambda_{\min}=\frac{\lambda}{mN}=5\,\mathrm{nm}.
\]
此处$m=1$，分辨本领100。首零半宽不是半高全宽；两者都随$1/N$变化，但系数不同。题设未给单缝宽度，因此只计算阵列峰的宽度和理想瑞利分辨率；该级能否被包络压低，还需单缝孔径信息。
\end{solution}
